\chapter{全文总结与展望}

\section{全文总结}
本文面向 BCI-IoT 场景的“敏感数据 + 实时性 + 资源受限”约束，设计并实现了 IoT\_BCI-Sudoku 安全通信增强协议。协议在身份认证握手与会话 AEAD 加密的基础上，引入 Sudoku 外观层（映射编码、随机填充与布局轮转），并提供抗重放、异常回落与可复现评测工具链。通过与纯 AEAD、DTLS、MQTT 与 CoAP/UDP 等基线协议的对比实验，本文给出在带宽开销、时延、内存与外观指标上的量化结果，并讨论了安全性、效率与抗侧写能力之间的工程权衡。

\section{后续工作展望}
结合当前实现与评测进度，后续可从以下方向进一步完善：
\begin{itemize}
  \item \textbf{端侧资源上限参数化}：为不同设备规格提供默认配置与上限保护（连接状态、缓冲区、回落策略阈值）。
  \item \textbf{更系统的安全验证}：补充 fuzz 测试覆盖握手/解码/状态机边界，并在 CI 中固定运行静态检查与竞态检测。
  \item \textbf{更贴近真实网络的极端场景}：增加可复现的高抖动/丢包/拥塞场景测试脚本与统计口径，用于评估鲁棒性与性能波动。
  \item \textbf{密钥生命周期管理}：支持证书吊销/轮换、会话密钥更新与多设备部署下的集中化管理。
  \item \textbf{进一步降低开销}：在不牺牲安全语义的前提下，持续优化内存分配、减少写调用次数，并探索更高效的 padding 策略。
\end{itemize}

